\chapter{Deterministic and Stochastic Correlation}

\section{Outcomes and Probabilities of Outcomes}

\subsection{Normal Distribution}
For mean $\mu \in \real$ and standard deviation 
$\sigma \in (0, \infty)$,
the {\bf normal distribution} has a 
probability density function (PDF) of 
\be
p(x) 
= 
\frac{1}{\sqrt{2 \pi \sigma^2}} 
\; \exp 
\left[ 
\frac{-(x - \mu)^2}{2 \sigma^2}
\right].
\ee
This PDF is valid since it is non-negative for values of
support $x \in \real$, and because the integral from 
$-\infty$ to $\infty$ is equal to unity.  
The {\bf cumulative distribution} 
function (CDF) is 
\be
P(x) = 
\int_{-\infty}^{x} p(\tau) \; d\tau = 
\half 
\left(
1 + 
\mbox{erf} \;
\left[\frac{x - \mu}{\sqrt{2 \sigma^2}}\right]
\right).
\ee
As $x \rightarrow -\infty$, $P \rightarrow 0$.
As $x \rightarrow \infty$, $P \rightarrow 1$.
The {\bf mean value} is 
\begin{align}
 \ex & = 
 \int_{-\infty}^{\infty} p(x) \cdot x \; dx \\
 & = 
 \int_{-\infty}^{\infty} 
\frac{x}{\sqrt{2 \pi \sigma^2}} 
\; \exp 
\left[ 
\frac{-(x - \mu)^2}{2 \sigma^2}
\right] \;dx \\
 & = \mu.
\end{align}
The {\bf variance} of $x$ is
\begin{align}
\var(X) 
 & = 
 \E\left[ (X - \ex)^2 \right] = \E[X^2] - (\ex)^2, \\
 & = 
 \int_{-\infty}^{\infty} p(x) \cdot x^2 \; dx - \mu^2,  \\
 & = 
 \int_{-\infty}^{\infty} 
\frac{x^2}{\sqrt{2 \pi \sigma^2}} 
\; \exp 
\left[ 
\frac{-(x - \mu)^2}{2 \sigma^2}
\right] \;dx - \mu^2,  \\
 & = \sigma^2.
\end{align}

\begin{Hremark}
For the normal distribution, the mean $\mu$ is also 
equal to the median and mode.  The skew and kurtosis are 
both zero.
\end{Hremark}

\begin{Hremark}
The 
normal distribution is also called the {\bf Gaussian
distribution}.  See 
Stahl\footnote{Stahl S. {\em The Evolution of the Normal
Distribution}, Mathematics Magazine, 79:2, April 2006.} 
for an account for 
how {\em normal} arose, and how this naming 
``... had the disadvantage of leading people to believe
that all other distributions of frequency are in one sense
or another {\em abnormal}.''
\end{Hremark}

\section{Expectation}

\begin{itemize}
\item Given a series of outcomes, expectation 
is a {\bf measure} of that series, which will 
be shaped by the probability of each of the outcomes.  
Expectation is also called the weighted average.
\item In the 
special case of equally likely probability for each outcome, 
the expectation will be the mean value, also called the 
simple average.
\end{itemize}

\subsection{Non-Equal Likelihood}

Let $X$ be a random variable with a finite number of 
{\bf outcomes} $x_i$
for $i=1\ldots n$ with probability $p_i$.  The {\bf expectation} 
of $X$, denoted as $\bar{x}$, is defined as
\be
\ex \defe \sum_{i=1}^n \; p_i \; x_i = \bar{x}.
\label{eq:simple_average}
\ee
Because the sum of all probabilities is identically 
equal to unity,
\be
\sum_{i=1}^n \; p_i = 1.0,
\ee
the expected
value is considered the 
{\bf weighted average}, where each $p_i$ is the 
{\bf weight}.

\begin{Hremark}
The weighted average is also 
known as the {\bf first moment}, which is 
analogous to the {\bf center of mass} of a rigid body.
The center of mass analogy applies to both the non-equal
likelihood case (non-uniform mass density) and equal
likelihood case (uniform mass density).
\end{Hremark}

\subsection{Equal Likelihood}
For the special case where each probability is equally likely, 
\be
p_1 = p_2 = \cdots = p_n = \frac{1}{n},
\label{eq:equal_probability}
\ee
the expectation of $\ex$ has the specialized form 
defined as $\mu$ 
\be
\ex \overset{(\ref{eq:equal_probability})}{=} 
\frac{1}{n} \sum_{i=1}^n \; x_i
\defe \mu,
\label{eq:simple_average}
\ee
and is referred to as the {\bf simple average}.

\begin{Hremark}
The expectation $\ex$ should not be confused 
with $\mu$, the simple average of $X$.  
The simple average (equal likelihood) is a special
case of the expectation (non-equal likelihood).
\end{Hremark}

\begin{Hremark}
When multiple variables are considered, \eg $X$ and $Y$, 
the simple averages are denoted 
$\mu_x$ and $\mu_y$ or equivalently, $\bar{x}$ and $\bar{y}$,
respectively.
\end{Hremark}

\begin{Hexample}
{
Consider a 6-sided dice, with faces 1 to 6.  Consider the dice
to be a {\em fair} dice, so that the probability of rolling
a 1 is the same as rolling a 2, and is the same as rolling any 
of the other remaining numbers 3 to 6.   Thus, the fair dice
has equal likelihood of any of the 6 outcomes, and 
Eq.~(\ref{eq:simple_average}) is used.
The expected value is calculated as
\be
\ex = \mu = \mu_x = \bar{x} = \frac{1}{6} \left(1 + 2 + 3 + 4 + 5 + 6 \right) = 3.5
\ee
}
\end{Hexample}
The 6-sided dice roll in the preceding example can be 
illustrated through numerical simulation, implemented in Python
below.  
The results from the Python script are shown in
Fig.~\ref{fig:dice_roll}.
\begin{minted}
[
frame=lines,
framesep=2mm,
baselinestretch=1.2,
bgcolor=hivory,
fontsize=\footnotesize,
linenos
]
{python}
#!/usr/bin/env python3
# dice_roll.py
import matplotlib.pyplot as plt
from matplotlib import rc
import numpy as np
from random import randint  # random integer
rc('font', **{'family': 'serif', 'serif': ['Computer Modern Roman']})
rc('text', usetex=True)

n_rolls = 40  # number of dice rolls
roll_number = np.arange(1, n_rolls + 1)
roll_history = np.zeros(n_rolls)  # each roll result recorded in time
average = np.zeros(n_rolls)  # initialize the average
# then update after each dice roll

for roll in roll_number:
    # roll a fair dice, values between 1 and 6, inclusive
    roll_outcome = randint(1, 6) 
    print(f'Dice roll number {roll} outcome is {roll_outcome}.')
    roll_history[roll-1] = int(roll_outcome)
    sum_after_roll = sum(i for i in roll_history)
    average[roll-1] = sum_after_roll / roll

fig = plt.figure(figsize=(8, 4))  # 8 inches wide, 4 inches tall

ax = fig.add_subplot(1, 1, 1)
ax.plot(roll_number, roll_history, 's', color='blue', label='roll outcome')
ax.plot(roll_number, average, color='red', marker='o', 
	alpha=0.5, label='roll average')
ax.grid()
ax.set_xlabel('roll number')
ax.set_ylabel('roll outcome and cumulative simple average')
ax.set_ylim(0, 7)
ax.legend(loc='upper right')
plt.show()
figure_name = 'dice_roll'
fig.savefig(figure_name + '.pdf', bbox_inches='tight')
\end{minted}

\begin{figure}[htb]
  \begin{center}
  %\begin{overpic}[width=\textwidth, grid, tics=10]{fig/1D_contact}
  %\begin{overpic}[width=\textwidth,natwidth=309,natheight=172]{fig/1D_contact.pdf}
  % \begin{overpic}[width=\textwidth]{fig/dice_roll.pdf}
  \begin{overpic}[width=\textwidth]{fig/dice_roll_final.pdf}
  %\put(17, 43.5){$\cA$}
 \end{overpic}
 \end{center}
 \caption{Numerical simulation of 40 dice rolls, showing value of 
 each roll and average for all rolls.  As the number of rolls 
 increases, the average converges to 3.5, the dice roll expected value.}
 \label{fig:dice_roll}
\end{figure}

The first moment analog to the center of mass 
can be illustrated through graphical inspection.  
Fig.~\ref{fig:dice_moment} shows the roll outcomes, 
$1, \ldots, 6$, along the $x$-axis and the probability of
any of the six outcomes, $p_i = 1/6$, normalized to unity 
(since all of the six probabilities are equal) on the $y$-axis.

The center of mass of the six blue squares is readily 
seen at the $x=3.5$ coordinate, indicated with
a red circle.  The center of mass is the simple average.  
This illustrates the equal probability special case, 
interpreted in the mechanical analog as the uniform density
case.

For the general case with non-equal probability, \eg
an unfair or weighted dice, the higher (or lower) probability of 
a particular outcome is analogous to an increase (or 
decrease) of local
density near the particular $x$-coordinate 
of that outcome, and hence as the 
non-uniform density case.

\begin{figure}[htb]
  \begin{center}
  \begin{overpic}[width=\textwidth]{fig/dice_moment.pdf}
 \end{overpic}
 \end{center}
 \caption{Graphical interpretation of the simple average as the first moment center of mass $\bar{x}$.}
 \label{fig:dice_moment}
\end{figure}

